\section{摆线齿轮}
\markboth{50. 摆线齿轮}{50. 摆线齿轮}

钟表和某些其它仪表中的齿轮，齿廓曲线不用圆的渐伸线，而是采用外摆线和内摆线。这样的齿轮叫做摆线齿轮。通过研究内摆线和外摆线的法线的性质，就能知道摆线齿轮啮合也能得到固定的传动比。

为此，我们把内摆线和外摆线的参数方程写成下面的统一形式（参考第 36 节）：
\begin{align*}
  x &= (R-mR)\cos m\theta + mR\cos(\theta-m\theta), \\
  y &= (R-mR)\sin m\theta - mR\sin(\theta-m\theta),
\end{align*}
其中当 $0<m<1$ 时为内摆线，$m>1$ 时为外摆线。用撇号表示对参数 $\theta$ 的导数，得到
\begin{align*}
  x' &= -m(R-mR)\sin m\theta - (1-m)mR\sin(\theta-m\theta) \\
     &= m(m-1)R\bigl[\sin m\theta + \sin(\theta-m\theta)\bigr], \\
  y' &= m(R-mR)\cos m\theta - (1-m)mR\cos(\theta-m\theta) \\
     &= m(m-1)R\bigl[-\cos m\theta + \cos(\theta-m\theta)\bigr].
\end{align*}
所以曲线在点 $M(x,y)$ 的切线 $MT$ 的斜率是
\begin{align*}
  k_1 = \frac{dy}{dx} &= \frac{y'\dd\theta}{x'\dd\theta} \\
                      &= \frac{-\cos m\theta + \cos(\theta-m\theta)}{\sin m\theta + \sin(\theta-m\theta)}.
\end{align*}

另一方面，如图 50.1，考虑圆周上的点 $N(R\cos m\theta,\, R\sin m\theta)$，直线 $MN$ 的斜率是
\[
  k_2 = \frac{y - R\sin m\theta}{x - R\cos m\theta}.
\]
把 $x$、$y$ 用参数表达式代入，得到
\begin{align*}
  k_2 &= \frac{[(R-mR)\sin m\theta - mR\sin(\theta-m\theta)] - R\sin m\theta}{[(R-mR)\cos m\theta + mR\cos(\theta-m\theta)] - R\cos m\theta} \\
      &= \frac{-mR[\sin m\theta + \sin(\theta-m\theta)]}{-mR[\cos m\theta - \cos(\theta-m\theta)]} \\
      &= -\frac{1}{k_1},
\end{align*}

\noindent
\begin{minipage}[c]{0.55\textwidth}
所以 $MN\perp MT$，即 $MN$ 是曲线的法线，这说明内摆线和外摆线在任意点 $M$ 的法线一定通过点 $N(R\cos m\theta,\, R\sin m\theta)$，这里 $\theta$ 是点 $M$ 的参数值。令
\[
  m\theta = \omega t,
\]
其中 $\omega$ 为常数，$t$ 为时间变量，那么随着 $t$ 的增大，点 $N$ 将会沿着圆周，以角速度 $\omega$ 绕 $O$ 点匀速旋转。
\end{minipage}%
\hfill
\begin{minipage}[c]{0.42\textwidth}
  \centering
  \begin{tikzpicture}[scale=0.90, >=Stealth, line cap=round, line join=round]
    \def\R{0.82}
    \def\m{2.6}
    \def\th{25} % m*th = 65 deg
    \pgfmathsetmacro{\mth}{\m*\th}
    \coordinate (O) at (0,0);
    \coordinate (B) at (\R, 0);
    \coordinate (N) at ({\R*cos(\mth)}, {\R*sin(\mth)});
    \coordinate (M) at ({\R*((1-\m)*cos(\mth) + \m*cos(\th-\mth))},
                        {\R*((1-\m)*sin(\mth) - \m*sin(\th-\mth))});
    % 切线方向 (dx, dy)
    \pgfmathsetmacro{\tx}{\m*(\m-1)*\R*(sin(\mth) + sin(\th-\mth))}
    \pgfmathsetmacro{\ty}{\m*(\m-1)*\R*(-cos(\mth) + cos(\th-\mth))}
    \pgfmathsetmacro{\tlen}{sqrt(\tx*\tx + \ty*\ty)}
    \pgfmathsetmacro{\ux}{\tx/\tlen}
    \pgfmathsetmacro{\uy}{\ty/\tlen}

    % 坐标轴
    \draw[->] (-1.10, 0) -- (1.75, 0) node[below, font=\scriptsize] {$x$};
    \draw[->] (0, -1.10) -- (0, 1.35) node[left, font=\scriptsize] {$y$};

    % 定圆 O
    \draw[thick] (O) circle (\R);

    % 外摆线弧 (theta from 0 to 46 deg)
    \draw[very thick, curveblue, domain=0:46, samples=60, variable=\t]
      plot ({\R*((1-\m)*cos(\m*\t) + \m*cos(\t-\m*\t))},
            {\R*((1-\m)*sin(\m*\t) - \m*sin(\t-\m*\t))});

    % 半径 ON 与法线 MN 延长线
    \draw (O) -- node[above left=-2pt, font=\scriptsize] {$R$} (N);
    \draw (N) -- ($(M)!-0.45!(N)$);

    % 切线 MT 穿过 M 与 M 处直角符号
    \coordinate (T1) at ($(M) - 0.65*(\ux, \uy)$);
    \coordinate (T2) at ($(M) + 0.55*(\ux, \uy)$);
    \draw[thick] (T1) -- (T2);
    \node[below, font=\scriptsize] at (T1) {$T$};
    \coordinate (Mn) at ($(M)!0.10cm!(N)$);
    \coordinate (Mt) at ($(M)+0.10*({\ux},{\uy})$);
    \draw[thin] (Mn) -- ($(Mn)+(Mt)-(M)$) -- (Mt);

    % 角度 m*theta
    \draw (0.24, 0) arc (0:\mth:0.24);
    \node[font=\scriptsize, fill=white, inner sep=0.3pt] at ({\mth/2}:0.43) {$m\theta$};

    % 标点与标签
    \fill (O) circle (1.2pt) node[below left=-1pt, font=\scriptsize] {$O$};
    \fill (B) circle (1.2pt) node[below right=0.5pt, xshift=-0.5pt, fill=white, inner sep=0.3pt, font=\scriptsize] {$B$};
    \fill (N) circle (1.2pt) node[above, font=\scriptsize] {$N$};
    \fill (M) circle (1.3pt) node[right, font=\scriptsize] {$M$};
  \end{tikzpicture}
  \captionof*{figure}{图\ 50.1}
\end{minipage}

\vspace{4pt}
如果与此同时，我们让图中的曲线和圆周一起绕 $O$ 点以角速度 $\omega$ 沿相反方向旋转，那么这两个相反的旋转就会互相抵消，使 $N$ 点始终停留在一个固定位置，即图中的定点 $B$ 处。因而，如果 $O$ 是一个齿轮的旋转中心，而把齿轮的齿廓曲线做成外摆线或内摆线，那么当这个齿轮以角速度 $\omega$ 旋转时，齿廓上动点 $M$ 处的法线总是通过定点 $B$。由此可设法选择一对摆线齿轮，并且把它们的中心间隔一个适当的距离，使得在传动过程中，这一对齿轮在啮合点处的公共法线总是通过连心线上的一个定点，因而根据第 49 节的讨论，这时一定得到固定的传动比。

从上面十分粗略的介绍中，已经可以看出，摆线齿轮在设计、制造、安装等方面的要求都比渐开线齿轮要高些，因而摆线齿轮的应用范围也不如渐开线齿轮广泛，在研究齿轮原理时，需要对曲线的知识作很多深入的研究。
